\documentclass[border=4pt]{standalone}
\usepackage{tikz}
\usetikzlibrary{arrows,positioning}

% Source-derived from the legacy open triangle 60 reversed declaration.
\makeatletter
\pgfarrowsdeclare{polar inclusion}{polar inclusion}
{
  \pgfutil@tempdima=.5pt
  \advance\pgfutil@tempdima by.25\pgflinewidth
  \pgfarrowsleftextend{-\pgflinewidth}
  \pgfutil@tempdimb=7.794\pgfutil@tempdima
  \advance\pgfutil@tempdimb by.5\pgflinewidth
  \pgfarrowsrightextend{\pgfutil@tempdimb}
}
{
  \pgfutil@tempdima=.5pt
  \advance\pgfutil@tempdima by.25\pgflinewidth
  \pgfsetmiterjoin
  \pgfpathmoveto{\pgfqpointpolar{30}{9\pgfutil@tempdima}}
  \pgfpathlineto{\pgfpointorigin}
  \pgfpathlineto{\pgfqpointpolar{-30}{9\pgfutil@tempdima}}
  \pgfpathclose
  \pgfusepathqstroke
}
\makeatother

\begin{document}
\begin{tikzpicture}[
  node distance=16mm and 25mm,
  space/.style={draw,circle,minimum size=10mm,inner sep=1pt}
]
  \node[space,fill=blue!8] (A) {$A$};
  \node[space,fill=green!10,right=of A] (B) {$B$};
  \node[space,fill=orange!10,below=of A] (C) {$C$};
  \node[space,fill=purple!8,right=of C] (D) {$D$};

  \draw[blue,line width=.4pt,-{polar inclusion}] (A) -- node[above] {$i$} (B);
  \draw[purple,line width=1pt,-{polar inclusion}] (C) -- node[below] {$j$} (D);
  \draw[red,line width=1.6pt,-{polar inclusion}] (A) -- node[left] {$k$} (C);
\end{tikzpicture}
\end{document}
